\documentclass[11pt]{article}

\usepackage[a4paper,margin=1in]{geometry}
\usepackage{amsmath,amssymb}
\usepackage{booktabs}
\usepackage{longtable}
\usepackage{array}
\usepackage{enumitem}
\usepackage{hyperref}
\usepackage{xcolor}

\hypersetup{
  colorlinks=true,
  linkcolor=blue!60!black,
  urlcolor=blue!60!black,
  citecolor=blue!60!black
}

\setlist[itemize]{leftmargin=1.4em}
\setlist[enumerate]{leftmargin=1.6em}
\renewcommand{\arraystretch}{1.15}

\title{Uniform Spanning Tree / LERW Random Environment Project Explainer}
\author{Alberto Rescigno}
\date{June 2026}

\begin{document}
\maketitle

\begin{abstract}
This note explains the whole current state of the numerical project on
loop-erased random walk (LERW) winding in fixed random environments. It covers
the mathematical motivation, the implemented models, the simulation pipeline,
the diagnostics used to detect drift, the focused numerical results up to
$L=256$, the supervisor feedback, and the next larger-run plan. The most
important current conclusion is that the original two-random-weight model
develops strong north-east drift after probability normalisation. The next
scientific direction is therefore to use locally balanced random environments,
where every lattice site receives its own independent sample and has no local
directional drift, and then test larger sizes and alternative distributions.
\end{abstract}

\tableofcontents
\newpage

\section{Executive Summary}

The project studies random walks on the square lattice box
\[
  B_L = [-L,L]^2 \cap \mathbb{Z}^2,
\]
started at the origin and stopped when the walk first hits the boundary. The
chronological loop erasure of the walk gives a loop-erased random walk path.
The main observable is the winding,
\[
  W_L = \#\{\text{left turns}\} - \#\{\text{right turns}\}.
\]

For the ordinary simple symmetric random walk, the expected qualitative
behaviour is
\[
  \operatorname{Var}(W_L) \sim C \log L.
\]
The random-environment question is whether a fixed site-dependent environment
can produce anomalous growth closer to
\[
  \operatorname{Var}(W_L) \sim C(\log L)^2.
\]

The initial implementation tested a two-random-weight edge model
\[
  w_N = 1,\qquad w_E = 1,\qquad w_S = u,\qquad w_W = v,
\]
with $u,v$ positive mean-one random variables. After normalisation into
transition probabilities, this model produces a strong effective north-east
drift for moderate and strong disorder. That drift makes the walk exit near the
north-east corner and suppresses winding variance. It is therefore not a clean
test of a centred random-environment winding question.

The latest supervisor feedback confirms the next direction:
\begin{itemize}
  \item use a more locally balanced random environment;
  \item ensure each site has a different random sample;
  \item run larger simulations;
  \item if the balanced Gamma model does not show clear $(\log L)^2$ behaviour
        at larger sizes, test other distribution choices.
\end{itemize}

The code now supports this plan through exact opposite-pair locally balanced
models: \texttt{gamma\_balanced}, \texttt{lognormal\_balanced},
\texttt{pareto\_balanced}, and \texttt{uniform\_balanced}. The older
four-independent-weight models \texttt{gamma4}, \texttt{lognormal4},
\texttt{pareto4}, and \texttt{uniform4} remain available as comparison models.

\section{Mathematical Background}

\subsection{Uniform Spanning Tree and LERW}

Wilson's algorithm samples a uniform spanning tree (UST) by repeatedly running
random walks and adding their loop erasures to the tree. Thus branches of the
UST are described by loop-erased random walk. Studying the geometry of LERW is
therefore a way of studying geometric properties of the UST.

In this project the numerical experiment is not sampling full spanning trees.
Instead, it focuses on one branch-like observable: run a random walk from the
origin until it exits a square box, erase loops chronologically, and measure
the winding of the resulting self-avoiding path.

\subsection{Winding Observable}

The square lattice directions are encoded cyclically as
\[
  E=0,\qquad N=1,\qquad W=2,\qquad S=3.
\]
For consecutive LERW steps, if the direction changes by $+1$ modulo $4$, this
is counted as a left turn. If it changes by $-1$ modulo $4$, this is counted as
a right turn. The winding is
\[
  W_L = L_L - R_L,
\]
where $L_L$ is the number of left turns and $R_L$ is the number of right turns
along the loop-erased path.

Immediate U-turns should not occur in a loop-erased nearest-neighbour path. The
code raises an error if one is encountered, because it would indicate a bug or
an invalid path.

\subsection{Scaling Question}

The core numerical question is whether the winding variance grows like
\[
  C \log L
\]
or whether a fixed random environment can produce a stronger law
\[
  C(\log L)^2.
\]
All numerical fits in this project are descriptive. They are not proofs of
asymptotic behaviour. With only finitely many sizes, especially sizes only up
to $L=256$, both $\log L$ and $(\log L)^2$ curves can sometimes look plausible.
The diagnostics must therefore be interpreted conservatively.

\section{Simulation Setup}

\subsection{Domain and Stopping Rule}

For each box size $L$, the walk starts at
\[
  X_0=(0,0)
\]
and evolves until the first time
\[
  \tau_L = \inf\{t \geq 0 : X_t \in \partial B_L\},
\]
where the boundary condition is implemented as
\[
  |x|=L \quad \text{or} \quad |y|=L.
\]
The raw path is
\[
  (X_0,X_1,\ldots,X_{\tau_L}).
\]
The code then chronologically erases loops to obtain the LERW path
\[
  \operatorname{LE}(X_0,\ldots,X_{\tau_L}).
\]

\subsection{Recorded Trial Data}

Each simulated trial records:
\begin{itemize}
  \item box size $L$;
  \item model name;
  \item distribution parameter, stored in the CSV column \texttt{k};
  \item sample index;
  \item random seed;
  \item raw walk length;
  \item LERW length;
  \item winding;
  \item boundary hit coordinates;
  \item number of random-environment sites sampled.
\end{itemize}

The main CSV schema is:
\begin{verbatim}
L,model,k,sample_index,seed,raw_steps,lerw_steps,winding,
hit_x,hit_y,environment_sites
\end{verbatim}

\subsection{Lazy Fixed Environments}

The random environment is sampled lazily. When the walk first visits a site,
the code samples that site's random weights and stores them in a dictionary.
Future visits to the same site reuse the same weights.

This is statistically equivalent to pre-sampling the entire box because the
site environments are independent. Laziness only avoids wasting time sampling
sites that the walk never visits.

This implementation directly addresses the supervisor's clarification that
each site should have a different sample.

\section{Implemented Models}

\subsection{Symmetric Baseline}

The calibration model is the ordinary simple symmetric random walk:
\[
  p_N=p_E=p_S=p_W=\frac14.
\]
This case is expected to show winding variance consistent with $C\log L$.

\subsection{Original Two-Random-Weight Gamma Model}

The original two-random-weight model uses
\[
  w_N=1,\qquad w_E=1,\qquad w_S=u,\qquad w_W=v,
\]
where
\[
  u,v \sim \Gamma(k,1/k).
\]
The scale $1/k$ makes the raw weights mean one:
\[
  \mathbb{E}[u]=\mathbb{E}[v]=1.
\]
Transition probabilities are obtained by normalising:
\[
  p_d(x)=\frac{w_d(x)}{w_N(x)+w_E(x)+w_S(x)+w_W(x)}.
\]

Although the raw random variables have mean one, this model is not locally
balanced after normalisation. Since north and east are fixed at weight $1$,
while south and west can be very small, the walk can acquire a strong effective
north-east drift.

\subsection{Exponential Two-Weight Model}

The \texttt{exponential} model has the same two-weight structure as the Gamma
model, but uses
\[
  u,v \sim \operatorname{Exp}(1).
\]
This is equivalent to the two-weight Gamma model at $k=1$, but it is kept as a
separate named model for clarity.

\subsection{Direction-Symmetric Gamma Model: \texttt{gamma4}}

The direction-symmetric Gamma model samples all four outgoing weights
independently:
\[
  w_N,w_E,w_S,w_W \stackrel{\text{iid}}{\sim} \Gamma(k,1/k).
\]
The four raw weights are then normalised to probabilities at that site.
Because all directions are treated symmetrically, this model is balanced in
distribution. However, at a single site it usually has non-zero local drift
because $w_N$ need not equal $w_S$ and $w_E$ need not equal $w_W$.

\subsection{Exact Locally Balanced Gamma Model: \texttt{gamma\_balanced}}

The main supervisor-directed model samples one vertical and one horizontal
weight at each site:
\[
  w_V,w_H \stackrel{\text{iid}}{\sim} \Gamma(k,1/k),
\]
then assigns
\[
  w_N=w_S=w_V,\qquad w_E=w_W=w_H.
\]
After normalisation, this gives
\[
  p_N=p_S,\qquad p_E=p_W,
\]
so every sampled site has zero local drift. Different sites still receive
different independent samples.

\subsection{Balanced Lognormal Models: \texttt{lognormal4} and
\texttt{lognormal\_balanced}}

The balanced lognormal model samples
\[
  w_N,w_E,w_S,w_W \stackrel{\text{iid}}{\sim}
  \operatorname{LogNormal}\left(-\frac{\sigma^2}{2},\sigma\right).
\]
The choice of mean parameter $-\sigma^2/2$ ensures
\[
  \mathbb{E}[w_d]=1.
\]
The model parameter is $\sigma$.
The exact local-balance version \texttt{lognormal\_balanced} uses the same law,
but samples only a vertical and a horizontal weight at each site and applies
them symmetrically to opposite directions.

\subsection{Balanced Pareto Models: \texttt{pareto4} and
\texttt{pareto\_balanced}}

The balanced Pareto model samples
\[
  w_d = \frac{\alpha-1}{\alpha}Y_d,
  \qquad
  Y_d \sim \operatorname{Pareto}(\alpha),
\]
independently for $d\in\{N,E,S,W\}$. The parameter must satisfy
\[
  \alpha>1
\]
so that the mean is finite. The prefactor makes the raw weights mean one.
This model gives a heavier-tailed alternative to Gamma and lognormal weights.
The exact local-balance version \texttt{pareto\_balanced} applies the same
opposite-pair construction as \texttt{gamma\_balanced}.

\subsection{Balanced Uniform Models: \texttt{uniform4} and
\texttt{uniform\_balanced}}

The balanced uniform model samples
\[
  w_N,w_E,w_S,w_W \stackrel{\text{iid}}{\sim} \operatorname{Uniform}(1-a,1+a),
\]
with
\[
  0<a<1.
\]
This gives a bounded-disorder comparison case. The exact local-balance version
\path{uniform_balanced} again uses opposite-pair vertical and horizontal
weights.

\section{Code Structure}

The current project folder contains the following main files.

\begin{description}[leftmargin=2.8cm,style=nextline]
  \item[\path{lerw_random_environment.py}]
  Core simulator. Defines the random environment, runs random walks to the
  boundary, performs loop erasure, computes winding, and writes trial CSVs.

  \item[\path{run_research_batch.py}]
  Batch runner. Runs many model/size/parameter configurations and supports
  per-configuration checkpoint files. It now supports exact local-balance
  configs through the \texttt{--configs} option.

  \item[\path{analyse_winding_results.py}]
  Command-line summary script. Reads one or more CSVs, groups by model,
  parameter, and size, then prints group means, variances, and simple scaling
  fits.

  \item[\path{make_figures.py}]
  Generates SVG figures, CSV tables, and a Markdown analysis dossier using
  only the Python standard library.

  \item[\path{test_lerw_random_environment.py}]
  Unit tests for loop erasure, winding, probability normalisation,
  fixed-in-time environments, and distinct site samples.

  \item[\path{README.md}]
  Practical project guide and current command reference.
\end{description}

\section{Existing Data and Outputs}

\subsection{Main Result CSVs}

Important result files include:
\begin{itemize}
  \item \path{results/research_batch_main.csv}: broad exploratory sweep.
  \item \path{results/research_batch_focused.csv}: focused run used for the
        main analysis below.
  \item \path{results/research_batch_focused_*.csv}: per-configuration
        checkpoints from the focused run.
  \item \path{results/smoke_balanced_distributions.csv}: small smoke test for
        the new balanced distributions.
\end{itemize}

\subsection{Analysis Output Folders}

\begin{itemize}
  \item \path{analysis_focused/}: focused figures, tables, and interpretation.
  \item \path{analysis_output/}: broader exploratory analysis.
  \item \path{analysis_smoke_balanced/}: smoke-test output for the new balanced
        distributions.
\end{itemize}

\section{Focused Numerical Run}

The focused run used:
\[
  L \in \{16,32,64,128,256\},
\]
with $700$ samples per model/size/parameter cell. The models were:
\begin{itemize}
  \item \texttt{symmetric};
  \item two-weight \texttt{gamma};
  \item balanced \texttt{gamma4}.
\end{itemize}
For \texttt{gamma} and \texttt{gamma4}, the tested parameters were
\[
  k \in \{20,1,0.5\}.
\]
The focused run contains $24{,}500$ trials.

\section{Focused Results: Symmetric Baseline}

The symmetric walk behaves as expected and calibrates the simulation.

\begin{center}
\begin{tabular}{rrrr}
\toprule
$L$ & $\operatorname{Var}(W)$ & mean raw steps & mean LERW steps \\
\midrule
16  & 1.894 & 296.9   & 40.3 \\
32  & 2.353 & 1239.1  & 96.2 \\
64  & 3.148 & 4861.2  & 231.8 \\
128 & 3.572 & 20045.7 & 559.1 \\
256 & 4.300 & 77271.2 & 1329.5 \\
\bottomrule
\end{tabular}
\end{center}

The descriptive fit against $\log L$ is
\[
  \operatorname{Var}(W) = 0.870 \log L - 0.565,
  \qquad R^2=0.992.
\]
The descriptive fit against $(\log L)^2$ is
\[
  \operatorname{Var}(W) = 0.104(\log L)^2 + 1.158,
  \qquad R^2=0.987.
\]
The log fit is slightly better, as expected for the baseline.

\section{Focused Results: Two-Weight Gamma Model}

\subsection{Drift Diagnostic}

At $L=256$, the mean boundary hit locations for the two-weight Gamma model are:

\begin{center}
\begin{tabular}{rrrr}
\toprule
$k$ & mean hit $x/L$ & mean hit $y/L$ & mean raw steps \\
\midrule
20  & 0.589 & 0.591 & 48969.1 \\
1   & 0.894 & 0.889 & 4731.9 \\
0.5 & 0.916 & 0.918 & 2891.9 \\
\bottomrule
\end{tabular}
\end{center}

This is the central modelling issue. For $k=1$ and $k=0.5$, the walk exits
very near the north-east corner. The drift is not a small numerical fluctuation.
It is large enough to change the qualitative behaviour of the raw walk.

\subsection{Winding Variance}

\begin{center}
\begin{tabular}{rrrrr}
\toprule
$k$ & $\operatorname{Var}(W)$ at $L=16$ &
$\operatorname{Var}(W)$ at $L=256$ &
$R^2$ for $\log L$ &
$R^2$ for $(\log L)^2$ \\
\midrule
20  & 1.973 & 3.735 & 0.952 & 0.908 \\
1   & 1.755 & 2.099 & 0.524 & 0.587 \\
0.5 & 1.358 & 1.548 & 0.950 & 0.916 \\
\bottomrule
\end{tabular}
\end{center}

The moderate and strong disorder cases have suppressed winding variance. This
is consistent with a drifted-walk regime: the walk exits quickly and does not
sample large-scale winding in a centred way.

\section{Focused Results: Balanced Gamma Model}

\subsection{Drift Diagnostic}

At $L=256$, the balanced \texttt{gamma4} model remains centred:

\begin{center}
\begin{tabular}{rrrr}
\toprule
$k$ & mean hit $x/L$ & mean hit $y/L$ & mean raw steps \\
\midrule
20  & 0.017 & -0.027 & 79695.7 \\
1   & 0.028 & 0.026  & 96802.5 \\
0.5 & 0.017 & -0.003 & 149512.6 \\
\bottomrule
\end{tabular}
\end{center}

This supports the interpretation that the four-weight model is a more suitable
candidate for a no-drift random-environment test.

\subsection{Winding Variance}

\begin{center}
\begin{tabular}{rrrrr}
\toprule
$k$ & $\operatorname{Var}(W)$ at $L=16$ &
$\operatorname{Var}(W)$ at $L=256$ &
$R^2$ for $\log L$ &
$R^2$ for $(\log L)^2$ \\
\midrule
20  & 2.034 & 4.104 & 0.963 & 0.938 \\
1   & 1.987 & 4.001 & 0.978 & 0.943 \\
0.5 & 1.965 & 4.184 & 0.976 & 0.946 \\
\bottomrule
\end{tabular}
\end{center}

Over $L\leq 256$, the balanced Gamma model still looks closer to logarithmic
growth than to clear $(\log L)^2$ growth. This does not disprove the conjectured
behaviour. It only says the currently tested size range is not enough to show
it convincingly.

\section{Interpretation of the Drift Problem}

The two-weight model uses mean-one raw weights, but mean-one raw weights do not
guarantee zero effective drift after normalisation. The key issue is structural:
north and east are fixed at weight $1$, while south and west are random. When
south and west weights are often small, north and east become relatively more
likely. This creates a systematic north-east preference.

The boundary-hit diagnostic is therefore essential. If the mean boundary hit
location is close to $(0,0)$ after scaling by $L$, the walk is plausibly
centred. If the mean hit is close to $(L,L)$, the walk is drifted and winding
variance should not be interpreted as evidence for the centred random
environment problem.

\section{Supervisor Feedback and Current Direction}

The latest feedback from Prof. Chhita confirms that the project should now use
a more locally balanced random environment. The important points are:
\begin{itemize}
  \item each site should have a different sample;
  \item larger simulations should be run;
  \item other distributions should be tested if the Gamma distribution does not
        show $(\log L)^2$ behaviour at larger sizes;
  \item the conjectured $(\log L)^2$ behaviour is of community interest, so
        identifying the right distributional setting is valuable.
\end{itemize}

The code changes made after this feedback implement exactly this workflow:
\begin{itemize}
  \item \texttt{gamma\_balanced} is the exact locally balanced Gamma model;
  \item \texttt{lognormal\_balanced}, \texttt{pareto\_balanced}, and
        \texttt{uniform\_balanced} are additional exact locally balanced
        distribution families;
  \item \texttt{gamma4}, \texttt{lognormal4}, \texttt{pareto4}, and
        \texttt{uniform4} remain available as four-independent-weight
        comparison models;
  \item the tests verify that probabilities are positive and normalised;
  \item the tests verify that different sites receive different random samples.
\end{itemize}

\section{Recommended Larger Simulation}

The recommended next large run is:

\begin{verbatim}
python3 run_research_batch.py \
  --sizes 16 32 64 128 256 512 \
  --samples 1000 \
  --seed 20260615 \
  --configs symmetric gamma_balanced:20 gamma_balanced:1 gamma_balanced:0.5 \
            lognormal_balanced:0.5 lognormal_balanced:1 \
            pareto_balanced:3 pareto_balanced:2 uniform_balanced:0.5 \
  --checkpoint-per-config \
  --output results/research_batch_balanced_large.csv
\end{verbatim}

The \texttt{--checkpoint-per-config} option is important. It writes a separate
CSV after each model/parameter configuration finishes, so partial results are
not lost if a long run is interrupted.

For a quick smoke test before the long run:

\begin{verbatim}
python3 run_research_batch.py \
  --sizes 16 32 \
  --samples 5 \
  --seed 20260615 \
  --configs symmetric gamma_balanced:1 lognormal_balanced:0.5 \
            pareto_balanced:2 uniform_balanced:0.5 \
  --output results/smoke_balanced_distributions.csv
\end{verbatim}

\section{How to Regenerate Figures and Tables}

For the focused old analysis:
\begin{verbatim}
python3 make_figures.py \
  results/research_batch_focused.csv \
  --output-dir analysis_focused
\end{verbatim}

For the new balanced large run, after it finishes:
\begin{verbatim}
python3 make_figures.py \
  results/research_batch_balanced_large.csv \
  --output-dir analysis_balanced_large
\end{verbatim}

The script writes:
\begin{itemize}
  \item \path{analysis_*/figures/*.svg};
  \item \path{analysis_*/tables/group_stats.csv};
  \item \path{analysis_*/tables/scaling_fits.csv};
  \item \path{analysis_*/analysis_dossier.md}.
\end{itemize}

\section{Figure Guide}

The most important generated figures are:
\begin{description}[leftmargin=3.8cm,style=nextline]
  \item[\path{variance_vs_logL_selected.svg}]
  Winding variance against $\log L$ for selected models.

  \item[\path{variance_vs_logL2_selected.svg}]
  Same data plotted against $(\log L)^2$.

  \item[\path{mean_hit_x_over_L.svg}]
  Drift diagnostic for the horizontal boundary hit coordinate.

  \item[\path{mean_hit_y_over_L.svg}]
  Drift diagnostic for the vertical boundary hit coordinate.

  \item[\path{boundary_hits_L128.svg}]
  Scatter plot of boundary hit locations at $L=128$ in the focused analysis.

  \item[\path{raw_walk_length.svg}]
  Mean raw walk length. Drifted models exit much faster.

  \item[\path{lerw_path_length.svg}]
  Mean loop-erased path length.

  \item[\path{environment_vector_fields.svg}]
  Local mean-step vector field visualisation.
\end{description}

\section{Testing and Verification}

The unit test command is:
\begin{verbatim}
python3 -m unittest -v
\end{verbatim}

The current test suite checks:
\begin{itemize}
  \item chronological loop erasure;
  \item winding sign convention;
  \item symmetric transition probabilities;
  \item positive and normalised probabilities for Gamma models;
  \item positive and normalised probabilities for lognormal, Pareto, and
        uniform balanced models;
  \item fixed-in-time environments at a site;
  \item different samples at different sites;
  \item small boundary-hitting trials.
\end{itemize}

At the time this explainer was written, the test suite passed:
\[
  12/12 \text{ tests passing.}
\]

\section{What Can Be Claimed Safely}

The following claims are currently defensible:
\begin{itemize}
  \item The symmetric baseline reproduces the expected qualitative logarithmic
        growth.
  \item The original two-weight Gamma model develops strong effective
        north-east drift after normalisation.
  \item Strong drift correlates with shorter raw walks and suppressed winding
        variance.
  \item The balanced four-weight Gamma model stays centred in the tested size
        range.
  \item Up to $L=256$, the balanced Gamma model does not show convincing
        $(\log L)^2$ growth.
  \item Larger sizes and alternative balanced distributions are the correct
        next numerical step.
\end{itemize}

The following claims should be avoided:
\begin{itemize}
  \item Do not claim the conjectured $(\log L)^2$ behaviour is false.
  \item Do not claim the random-environment problem is solved.
  \item Do not call the two-weight model drift-free just because the raw
        weights have mean one.
  \item Do not treat one balanced distribution as decisive without checking
        larger sizes and other distribution families.
\end{itemize}

\section{Potential Paper Structure}

A future write-up could be structured as follows:
\begin{enumerate}
  \item Introduction: UST, LERW, winding, and motivation for random
        environments.
  \item Model definitions: symmetric walk, two-weight model, balanced
        four-weight models.
  \item Numerical method: box, stopping rule, loop erasure, winding convention,
        seeding, sample sizes.
  \item Baseline validation: symmetric $C\log L$ behaviour.
  \item Drift diagnostic: two-weight Gamma model and north-east exit bias.
  \item Balanced environments: Gamma and alternative distributions.
  \item Scaling diagnostics: fits to $\log L$ and $(\log L)^2$.
  \item Discussion: finite-size limitations, drift versus disorder, and future
        larger runs.
  \item Conclusion: what is established and what remains open.
\end{enumerate}

\section{Open Questions}

The main open scientific questions are:
\begin{itemize}
  \item Does any locally balanced fixed random environment show convincing
        $(\log L)^2$ winding variance at larger $L$?
  \item If Gamma weights do not show the effect, do heavier-tailed exact
        local-balance models such as \texttt{pareto\_balanced} produce it?
  \item How large must $L$ be before the asymptotic regime becomes visible?
  \item How many samples are needed for stable variance estimates at $L=512$
        or larger?
  \item Should future diagnostics include bootstrap confidence intervals or
        error bars for the variance estimates?
\end{itemize}

\section{Immediate Next Steps}

\begin{enumerate}
  \item Run the documented smoke test if needed.
  \item Launch the larger balanced batch with checkpointing.
  \item Regenerate figures and tables for the new large CSV.
  \item Check boundary-hit diagnostics first. Any model with strong drift
        should be treated as a drifted model, not a centred random environment.
  \item Compare variance fits against $\log L$ and $(\log L)^2$.
  \item Prioritise models that are both centred and show possible upward
        curvature relative to $\log L$.
  \item Add error bars before making stronger claims.
\end{enumerate}

\appendix

\section{Command Reference}

\subsection{Run Tests}
\begin{verbatim}
python3 -m unittest -v
\end{verbatim}

\subsection{Single-Model Debug Run}
\begin{verbatim}
python3 lerw_random_environment.py \
  --sizes 16 32 64 \
  --samples 50 \
  --model gamma_balanced \
  --k 1 \
  --output results/debug_gamma_balanced_k1.csv
\end{verbatim}

\subsection{Previous Focused Batch}
\begin{verbatim}
python3 run_research_batch.py \
  --sizes 16 32 64 128 256 \
  --samples 700 \
  --seed 20260613 \
  --models symmetric gamma gamma4 \
  --k-values 20 1 0.5 \
  --checkpoint-per-config \
  --output results/research_batch_focused.csv
\end{verbatim}

\subsection{New Balanced Batch}
\begin{verbatim}
python3 run_research_batch.py \
  --sizes 16 32 64 128 256 512 \
  --samples 1000 \
  --seed 20260615 \
  --configs symmetric gamma_balanced:20 gamma_balanced:1 gamma_balanced:0.5 \
            lognormal_balanced:0.5 lognormal_balanced:1 \
            pareto_balanced:3 pareto_balanced:2 uniform_balanced:0.5 \
  --checkpoint-per-config \
  --output results/research_batch_balanced_large.csv
\end{verbatim}

\subsection{Analyse a CSV in the Terminal}
\begin{verbatim}
python3 analyse_winding_results.py results/research_batch_focused.csv
\end{verbatim}

\subsection{Generate Figures}
\begin{verbatim}
python3 make_figures.py \
  results/research_batch_focused.csv \
  --output-dir analysis_focused
\end{verbatim}

\section{Glossary}

\begin{description}[leftmargin=3.0cm,style=nextline]
  \item[LERW]
  Loop-erased random walk. It is obtained from a random walk by erasing loops in
  chronological order.

  \item[UST]
  Uniform spanning tree. Wilson's algorithm constructs it using LERW.

  \item[Fixed environment]
  Random transition probabilities are sampled once per site and then kept fixed
  throughout the walk.

  \item[Locally balanced]
  All four outgoing directions at a site are sampled symmetrically from the
  same distribution, so there is no built-in preferred direction in law.

  \item[Drift diagnostic]
  A check based on the mean boundary hit location. Values of mean hit $x/L$ and
  $y/L$ close to zero indicate a centred walk. Values close to one indicate
  north-east drift.

  \item[Disorder strength]
  Informal term for how variable the random weights are. Smaller Gamma $k$,
  larger lognormal $\sigma$, and smaller Pareto $\alpha$ generally correspond
  to stronger disorder.
\end{description}

\end{document}
